\section{Introduction and Scope}
\label{sec:intro}

Corrosion in marine and offshore systems is a first-order economic and safety problem. The global cost of corrosion has been estimated at roughly USD~2.5~trillion per year, on the order of 3--4\% of world gross domestic product, with marine infrastructure, shipping, and offshore energy accounting for a substantial share \cite{koch2016}. Seawater is among the most aggressive natural electrolytes: a near-neutral solution containing approximately 0.5--0.6~mol/L of chloride, dissolved oxygen, biological activity, and wide swings in temperature, flow velocity, and aeration between the atmospheric, splash, tidal, and fully immersed zones \cite{nass2025marine}.

Conventional marine alloys each carry well-known limitations. Austenitic 316L stainless steel relies on a chromium-rich passive film but remains susceptible to pitting, crevice corrosion, and stress corrosion cracking (SCC) in warm, stagnant, or creviced chloride service \cite{lpbf_316l_scc2022}. Nickel--aluminum bronzes resist biofouling reasonably well but suffer dealuminification and cavitation damage; copper--nickel alloys trade strength for antifouling behavior; duplex stainless steels and titanium alloys perform well but are expensive and difficult to fabricate into complex or large geometries \cite{nass2025marine}. A material family that combines strength, toughness, and robust passivity in chloride media---and that can be formed into complex, repairable components---would therefore address a genuine gap.

High-entropy alloys (HEAs), defined by five or more principal elements in near-equimolar or deliberately non-equimolar ratios \cite{yeh2004,cantor2004}, are credible candidates. Their high configurational entropy, severe lattice distortion, sluggish diffusion, and the so-called ``cocktail'' effect together stabilize simple solid-solution phases and, in Cr-bearing compositions, support protective passive films \cite{miracle2017,george2019,murty2019}. Reviews of HEA corrosion behavior identify compositions---particularly FCC CoCrFeMnNi-type and carefully balanced Al-containing eutectic alloys---whose pitting resistance in chloride solutions rivals or exceeds 316L \cite{qiu2017,qiu2019}.

Additive manufacturing (AM) completes the picture for three reasons. \emph{First}, AM's rapid solidification ($\sim10^{3}$--$10^{6}$~K/s in powder-bed processes) suppresses the coarse segregation typical of cast HEAs, and recent evidence shows that this homogeneity can translate directly into better chloride corrosion performance \cite{slm_ehea_2025,han2025}. \emph{Second}, AM enables geometries that casting cannot deliver economically: conformally cooled heat exchangers, topologically optimized brackets, and integrated manifolds for seawater systems \cite{cagirici2024}. \emph{Third}, directed energy deposition (DED) and wire arc additive manufacturing (WAAM) make HEAs practical as repair and cladding materials for large marine components, avoiding the need to build entire structures from expensive alloys \cite{jiang2019clad,raultaiwade2021}.

Existing reviews have treated AM of HEAs comprehensively from the processing and mechanical-property viewpoints \cite{cagirici2024,han2025,abdullah2024rhea}, and marine-corrosion texts mention HEAs as an emerging material class \cite{nass2025marine}, but no synthesis to date is organized around \emph{marine deployment}: which compositions survive real seawater service, which AM routes suit which marine components, and what qualification evidence is missing. This review targets that gap. Its scope is: (i) HEA fundamentals relevant to marine service (Section~\ref{sec:fundamentals}); (ii) AM processes and feedstock for HEAs (Section~\ref{sec:am}); (iii) microstructure and mechanical performance of AM HEAs (Section~\ref{sec:micro}); (iv) corrosion and electrochemical behavior in saline and seawater environments (Section~\ref{sec:corrosion}); (v) candidate marine applications and qualification pathways (Section~\ref{sec:applications}); and (vi) challenges and a deployment roadmap (Sections~\ref{sec:challenges}--\ref{sec:conclusion}). Table and figure sets are collected in Section~\ref{sec:tabfig}.

\subsection{Review methodology and literature coverage}
\label{sec:method}

To formalize coverage and limit selection bias, the literature search followed a PRISMA-informed protocol \cite{prisma2020}, executed on 8 October 2026 (Fig.~\ref{fig:prisma}). \emph{Sources}: 26 structured queries run through a web-search interface surfacing indexed sources (ScienceDirect, SpringerLink, MDPI, Frontiers, Nature portfolio, Wiley, IOP, ASTM, OSTI, NRC, ResearchGate/Semantic Scholar records), plus reference-list snowballing from six AM-HEA and marine-corrosion reviews. \emph{Keywords}: built from three facets --- alloy family (``high-entropy alloy'', ``CoCrFeMnNi'', ``eutectic high-entropy alloy'', ``refractory HEA'', ``medium-entropy alloy''), process (``laser powder bed fusion'', ``SLM'', ``directed energy deposition'', ``wire arc additive manufacturing'', ``binder jetting'', ``laser cladding''), and marine/corrosion service (``seawater'', ``chloride'', ``pitting potential'', ``crevice'', ``stress corrosion cracking'', ``microbiologically influenced corrosion'', ``cathodic protection'', ``marine'') --- combined with benchmark terms (``316L'', ``NiAl bronze''). No date restriction was imposed; records span 2004--2026 with emphasis on 2019--2026. \emph{Inclusion criteria}: peer-reviewed primary studies or verifiable technical reports on AM HEAs (or AM stainless analogs where defect--corrosion links were needed) reporting microstructure, mechanical, electrochemical, corrosion, MIC, or standards content relevant to marine service. \emph{Exclusion criteria}: non-English records; records without verifiable bibliographic identity; duplicate studies superseded by stronger evidence; off-topic HEA applications. Counts as recorded: $\approx$105 records surfaced by queries, $\approx$35 by snowballing, $\approx$120 unique records screened, $\approx$59 excluded, 61 included (see \texttt{references.bib}). \emph{Declared limitation}: subscription databases (Scopus/Web of Science) were not accessible from the review environment, so coverage is structured and transparent but not exhaustive; gray-literature records were included only when bibliographically verifiable, and unverifiable details are flagged as such in the Data Appendix.
